\section{Validation Approach}

\subsection{Objectives}
The validation approach verifies that the historian correctly receives, stores, and exposes PTP-timestamped measurements from the DAQ nodes. Three properties are of primary interest:

\begin{itemize}[noitemsep]
    \item \textbf{Correctness}: data written to InfluxDB matches the values published by the STM32 nodes
    \item \textbf{Timestamp fidelity}: the PTP timestamp from the STM32 is preserved end-to-end with no substitution or modification
    \item \textbf{Cross-node alignment}: points from two synchronised nodes at the same physical instant share consistent timestamps within the expected PTP accuracy bound
\end{itemize}

\subsection{Component-Level Tests}
Before end-to-end validation, each pipeline stage is verified in isolation.

\begin{table}[H]
\centering
\caption{Component-Level Validation Tests}
\begin{tabular}{p{0.2\linewidth}p{0.35\linewidth}p{0.35\linewidth}}
\hline
\textbf{Test} & \textbf{Method} & \textbf{Pass condition} \\
\hline
Telegraf MQTT subscription & Publish a known burst with \texttt{mosquitto\_pub}; inspect Telegraf logs & Payload parsed, point written to \texttt{telemetry} bucket \\
Timestamp preservation & Publish a burst with a known \texttt{ts}; query InfluxDB and compare & Stored \texttt{\_time} matches \texttt{ts} to nanosecond precision \\
Tag extraction & Publish to a known topic; query tags on the stored point & \texttt{node\_id}, \texttt{subsystem}, and \texttt{metric} match the topic segments \\
Burst flattening & Publish a 20-sample burst; query the stored point & Fields \texttt{v\_0} to \texttt{v\_19} hold the published values in order; \texttt{sync} and \texttt{dt\_ns} present \\
\hline
\end{tabular}
\end{table}

\subsection{End-to-End Validation}
End-to-end validation uses the two-node testbed with both STM32 nodes publishing live measurements.

\textbf{Timestamp fidelity test:} A known electrical event (e.g.\ a step load change) is applied. The time of the event as recorded in the historian for each node, reconstructed from \texttt{\_time} and \texttt{dt\_ns}, is compared. The difference should be within the expected PTP synchronisation accuracy. A large discrepancy indicates either a timestamp substitution error in the pipeline or a PTP synchronisation failure. Because each node publishes only one 20~ms window every 5~s, the event must fall within a published window on both nodes; the full-rate binary stream removes this constraint.

\textbf{Data completeness test:} Both nodes publish for a fixed duration. The number of points in the historian for each channel is compared against the expected count of one burst per channel every 5~s. Any gap is investigated against Telegraf and Mosquitto logs and the node's UART counters.

\textbf{PTP lock test:} The \texttt{sync} field is queried over a sustained period. It should remain 1 throughout normal operation once the servo has locked.

\subsection{Instrumentation}
Ingestion health is monitored through Telegraf's container logs (\texttt{docker compose logs telegraf}) and, on the node side, the acquisition and MQTT counters printed once per second on each node's debug UART.

\paragraph{Deferred} Telegraf's \texttt{internal} input plugin, writing pipeline metrics (points received, write errors, flush latency) to a separate \texttt{telegraf\_internal} bucket, was designed but not enabled.
